\section{Convert a Min Heap to a Max Heap}
\label{sec:heap-min-to-max}

\problemstatement{Given an array that currently satisfies the min-heap
property, rearrange it in place so that it satisfies the max-heap property
instead.}

\begin{patternbox}
The keyword here is \emph{``convert,'' not ``rebuild from a sorted
list''} -- the existing min-heap order is, in fact, almost entirely
irrelevant to solving this efficiently; the fastest approach simply
\textbf{re-heapifies the array from scratch as a max-heap}, ignoring its
current (min-heap) structure. This problem is the natural place to
introduce the classic $O(n)$ \textbf{build-heap} algorithm, which is
faster than the $O(n \log n)$ you would get by popping every element from
the min-heap and pushing it into a fresh max-heap one at a time.
\end{patternbox}

\subsection*{Understanding the problem carefully}

The simplest correct approach: treat the existing array's \emph{values}
as an unordered list of $n$ numbers (discard the fact that they currently
satisfy a min-heap arrangement -- it buys us nothing for the new task),
and re-heapify in place into a max-heap. The key efficiency insight is
that \emph{any} array can be turned into a valid heap in $O(n)$ time, not
$O(n \log n)$, by sifting down starting from the \textbf{last internal
node} and working backward to the root.

\begin{ideabox}[Why Starting from the Last Internal Node Works]
A leaf node trivially satisfies the heap property on its own. If we
process internal nodes in \emph{reverse} level order (bottom-most,
right-most internal node first, ending at the root), then by the time we
sift-down a given node $i$, both of its children's subtrees have
\emph{already} been fixed into valid max-heaps (since they were processed
earlier, being lower in the tree). Sifting node $i$ down then only needs
to resolve a violation between $i$ and the roots of two already-valid
subtrees -- exactly the situation sift-down is designed to fix -- so a
single bottom-up pass, sifting each internal node down once, is enough to
turn the entire array into a valid heap.
\end{ideabox}

\subsection*{Why this build-heap process is $O(n)$, not $O(n \log n)$}

Each sift-down costs time proportional to the \emph{height of the subtree
rooted at that node}, not the height of the whole tree. Most nodes are
near the bottom of the tree, where subtree heights are tiny; only the
handful of nodes near the root have a large subtree height. Formally, the
total work is
\[
  \sum_{h=0}^{\log n} \frac{n}{2^{h+1}} \cdot h = O(n),
\]
since the number of nodes at height $h$ from the bottom is roughly
$n/2^{h+1}$, and this geometric-times-linear sum converges to $O(n)$
rather than the naive (and incorrect) $O(n \log n)$ upper bound you would
get by assuming every node costs the full tree height.

\subsection*{Algorithm}

\begin{enumerate}
  \item The last internal node (one with at least one child) is at index
        $\lfloor n/2 \rfloor - 1$.
  \item For $i$ from $\lfloor n/2 \rfloor - 1$ down to $0$: sift-down
        node $i$, using max-heap comparisons (swap with the \emph{larger}
        child whenever a child exceeds the current node).
  \item The array now satisfies the max-heap property.
\end{enumerate}

\begin{algorithm}[H]
\caption{Build a Max-Heap In Place (Heapify)}
\begin{algorithmic}[1]
\Procedure{BuildMaxHeap}{\textit{arr}}
  \For{$i \gets \lfloor n/2 \rfloor - 1$ \textbf{down to} $0$}
    \State \Call{SiftDownMax}{\textit{arr}, $i$, $n$}
  \EndFor
\EndProcedure
\Procedure{SiftDownMax}{\textit{arr}, $i$, $n$}
  \While{\textbf{true}}
    \State $\textit{largest} \gets i$, $\textit{left} \gets 2i+1$, $\textit{right} \gets 2i+2$
    \If{$\textit{left} < n$ \textbf{and} $\textit{arr}[\textit{left}] > \textit{arr}[\textit{largest}]$}
      \State $\textit{largest} \gets \textit{left}$
    \EndIf
    \If{$\textit{right} < n$ \textbf{and} $\textit{arr}[\textit{right}] > \textit{arr}[\textit{largest}]$}
      \State $\textit{largest} \gets \textit{right}$
    \EndIf
    \If{$\textit{largest} = i$} \textbf{break} \EndIf
    \State swap $\textit{arr}[i]$ and $\textit{arr}[\textit{largest}]$; $i \gets \textit{largest}$
  \EndWhile
\EndProcedure
\end{algorithmic}
\end{algorithm}

\subsection*{Dry run}

\dryrun{Take the min-heap $\textit{arr} = [1,3,5,4,6,7,8]$ ($n=7$).}

Last internal node: $\lfloor 7/2 \rfloor - 1 = 2$.

\begin{itemize}
  \item $i=2$ (value $5$): children at $5,6$ are $7,8$. Largest child is
        $8$ at index $6$. $8>5$: swap. $\textit{arr}=[1,3,8,4,6,7,5]$.
        Index $6$ is a leaf; stop.
  \item $i=1$ (value $3$): children at $3,4$ are $4,6$. Largest is $6$ at
        index $4$. $6>3$: swap. $\textit{arr}=[1,6,8,4,3,7,5]$. Index $4$
        is a leaf; stop.
  \item $i=0$ (value $1$): children at $1,2$ are $6,8$. Largest is $8$ at
        index $2$. $8>1$: swap. $\textit{arr}=[8,6,1,4,3,7,5]$. Now at
        index $2$ (value $1$): children at $5,6$ are $7,5$. Largest is
        $7$ at index $5$. $7>1$: swap.
        $\textit{arr}=[8,6,7,4,3,1,5]$. Index $5$ is a leaf; stop.
\end{itemize}

The final array $[8,6,7,4,3,1,5]$ satisfies the max-heap property (the
root $8$ is the maximum, and every parent exceeds its children).

\subsection*{C++ implementation}

\begin{lstlisting}
#include <bits/stdc++.h>
using namespace std;

void siftDownMax(vector<int>& arr, int i, int n) {
    while (true) {
        int largest = i, left = 2 * i + 1, right = 2 * i + 2;
        if (left < n && arr[left] > arr[largest]) largest = left;
        if (right < n && arr[right] > arr[largest]) largest = right;
        if (largest == i) break;
        swap(arr[i], arr[largest]);
        i = largest;
    }
}

void minHeapToMaxHeap(vector<int>& arr) {
    int n = arr.size();
    for (int i = n / 2 - 1; i >= 0; i--) {
        siftDownMax(arr, i, n);
    }
}

int main() {
    vector<int> arr = {1, 3, 5, 4, 6, 7, 8}; // valid min-heap
    minHeapToMaxHeap(arr);

    cout << "Max-heap: ";
    for (int x : arr) cout << x << " ";
    cout << endl;
    return 0;
}
\end{lstlisting}

\begin{complexitybox}
\textbf{Time Complexity.} The build-heap process visits $O(n/2)$ internal
nodes with total sift-down work summing to
\[
  O(n)
\]
across the whole array (not $O(n \log n)$), by the height-weighted
geometric sum argument above.
\textbf{Space Complexity.} $O(1)$ auxiliary space (the conversion is done
in place).
\end{complexitybox}

\begin{takeawaybox}
Re-heapifying an entire array from scratch is $O(n)$, not $O(n \log n)$ --
a fact that is easy to misremember because a \emph{single} sift-down is
$O(\log n)$, tempting you to multiply by $n$ insertions. The $O(n)$ bound
only holds when you build the whole heap at once via bottom-up sift-down,
not when you insert $n$ elements one at a time via repeated sift-up.
\textsc{std::make\_heap} in C++ uses exactly this $O(n)$ bottom-up
technique internally.
\end{takeawaybox}
